%% ma: L11-B16-0096
%% lop: 11
%% chuong: 5
%% bai: 16
%% dang: 11.16.1
%% loai: TL
%% muc: VDC
%% dap_an: "-"
%% nguon: "Ảnh giáo viên gửi 25/11/2026: Giới hạn hàm số dạng 0/0 (3 bài, có ĐS), Bài 3 câu 4"
%% kien_thuc: "Lớp 11 Bài 16 (giới hạn của hàm số); hằng đẳng thức x^k-1=(x-1)(x^(k-1)+...+1); tổng 1+2+...+n (Lớp 11 Bài 6 cấp số cộng)"
%% trang_thai: chinh-thuc
%% ghi_chu: "Đề gốc không ghi điều kiện của n, đã thêm n là số nguyên dương."
\begin{ex}
Tính giới hạn $A=\lim\limits_{x\to1}\dfrac{x^{n+1}-(n+1)x+n}{(x-1)^2}\ (n\text{ là số nguyên dương})$.
\loigiai{Như trên với $n+1$ thay cho $n$: $x^{n+1}-(n+1)x+n=(x-1)^2\left[(x^{n-1}+\dots+1)+\dots+1\right]$. Vậy $A=n+(n-1)+\dots+1=\dfrac{n(n+1)}2$.}
\end{ex}
